% Copyright 2018  李文威 (Wen-Wei Li).
% Permission is granted to copy, distribute and/or modify this
% document under the terms of the Creative Commons
% Attribution 4.0 International (CC BY 4.0)
% http://creativecommons.org/licenses/by/4.0/

% 目的: 字体相关设置, 呼叫相关宏包.
% 将由 AJbook.cls 引入
% 必须提供 \kaishu, \songti, \heiti, \thmheiti, \fangsong 几种字型切换命令, 在文档类中使用.
\ProvidesFile{font-setup-local.tex}[2026/09/20 Project font files]

% 设置 xeCJK 字体及中文数字
%\setmainfont{TeX Gyre Pagella}	% 设置西文衬线字体
%\setsansfont{TeX Gyre Heros}[	% 设置西文无衬线字体
\setsansfont{texgyreheros}[     % 设置西文无衬线字体
Extension=.otf,
UprightFont=*-regular,
BoldFont=*-bold,
ItalicFont=*-italic,
BoldItalicFont=*-bolditalic]

% 使用用户提供的项目字体；从 notes/ 编译，源码包保留同级 font/。
% 宋体正文；黑体用于标题和粗体；楷体用于中文斜体；仿宋用于定理正文。
\setCJKmainfont[Path=../font/,
  BoldFont=FZHTK.TTF, ItalicFont=FZKTK.TTF,
  BoldItalicFont=FZHTK.TTF]{FZSSK.TTF}
\setCJKsansfont[Path=../font/,
  BoldFont=FZHTK.TTF, ItalicFont=FZHTK.TTF,
  BoldItalicFont=FZHTK.TTF]{FZHTK.TTF}
\setCJKmonofont[Path=../font/,
  BoldFont=FZHTK.TTF, ItalicFont=FZHTK.TTF,
  BoldItalicFont=FZHTK.TTF]{FZHTK.TTF}

\setCJKfamilyfont{song}[Path=../font/,
  BoldFont=FZHTK.TTF, ItalicFont=FZKTK.TTF,
  BoldItalicFont=FZHTK.TTF]{FZSSK.TTF}
\setCJKfamilyfont{kai}[Path=../font/,
  BoldFont=FZHTK.TTF, ItalicFont=FZKTK.TTF,
  BoldItalicFont=FZHTK.TTF]{FZKTK.TTF}
\setCJKfamilyfont{fangsong}[Path=../font/,
  BoldFont=FZHTK.TTF, ItalicFont=FZKTK.TTF,
  BoldItalicFont=FZHTK.TTF]{FZFSK.TTF}
\setCJKfamilyfont{hei}[Path=../font/,
  BoldFont=FZHTK.TTF, ItalicFont=FZHTK.TTF,
  BoldItalicFont=FZHTK.TTF]{FZHTK.TTF}
\setCJKfamilyfont{hei2}[Path=../font/,BoldFont=FZHTK.TTF]{FZHTK.TTF}
\setCJKfamilyfont{sectionfont}[Path=../font/,BoldFont=FZHTK.TTF]{FZHTK.TTF}
\setCJKfamilyfont{pffont}[Path=../font/,BoldFont=FZHTK.TTF]{FZHTK.TTF}
\setCJKfamilyfont{emfont}[Path=../font/,BoldFont=FZHTK.TTF]{FZHTK.TTF}

\defaultfontfeatures{Ligatures=TeX} 
\XeTeXlinebreaklocale "zh"
\XeTeXlinebreakskip = 0pt plus 1pt minus 0.1pt

% 以下设置字体相关命令, 用于定理等环境中.
\newcommand\kaishu{\CJKfamily{kai}} % 楷体
\newcommand\songti{\CJKfamily{song}} % 宋体
\newcommand\heiti{\CJKfamily{hei}}	% 黑体
\newcommand\thmheiti{\CJKfamily{hei2}}	% 用于定理名称的黑体
\newcommand\fangsong{\CJKfamily{fangsong}} % 仿宋
\renewcommand{\em}{\bfseries\CJKfamily{emfont}} % 强调